% Einheit 3 von 5 – Dreiecke nachweisen (bank/punkte-und-strecken-im-koordinatensystem/e3.jsonl, zusammenbau v0.1)
%% TODO Zeitmarke Einheit 3: Marken-Zeile nicht lesbar
%% TODO Zweigzeile Teil 1: was hier gelernt wird (Satz in Schülersprache aus dem Einheitstitel) – Einheit 3
\einheitenkopf[e3]{Einheit 3 von 5 · Dreiecke nachweisen}
\zweigzeile{Abitur GK}
\setcounter{aufgabe}{14}

%% TODO Titel: Ich-kann-Satz fehlt in der Bank (Platzhalter: Kettenname) – Nr. 15
%% TODO Anweisung: ein Satz über der ersten Teilaufgabe fehlt in der Bank – Nr. 15
\begin{aufgabe}{Dreiecke nachweisen}
%% TODO \rechenplatz{n}: Schrittzahl des Grundfalls fehlt in der Bank (nur bei fester Schreibform, 2.3 b)
\begin{teile}
\teil Parallelogramm $ABCD$: Welcher Vektor ist gleich $\overrightarrow{AB}$? \\ \kreuz{$\overrightarrow{CD}$} \\ \kreuz{$\overrightarrow{DC}$} \\ \kreuz{$\overrightarrow{BC}$}
\teil Zeige: Das Dreieck $A(2 | 0 | 1)$, $B(6 | 2 | 1)$, $C(4 | 1 | 5)$ ist gleichschenklig. Nenne Basis und Schenkel. \hfill (Abitur 2023 GK)
\teil Prüfe, ob das Dreieck $A(4 | 2 | 2)$, $B(2 | 4 | 2)$, $C(2 | 2 | 4)$ gleichseitig ist. \hfill (Abitur 2023 GK)
\teil $A(1 | 2 | 0)$, $B(5 | 6 | 0)$ und $C(3 | 4 | p)$ mit $p > 0$. Zeige: Das Dreieck $ABC$ ist für jedes $p$ gleichschenklig mit der Basis $AB$. \hfill (Abitur 2022 GK)
\teil Pyramiden mit $A(0 | 0 | 0)$, $B(7 | 0 | 0)$, $C(0 | 7 | 0)$ und der Spitze $D_p(0 | 0 | p)$, $p > 0$. Begründe ohne Rechnung: Das Dreieck $BCD_p$ ist gleichschenklig. \hfill (Abitur 2022 LK)
\teil Grundfläche $A(6 | 1 | 0)$, $B(1 | 4 | 0)$, $C(1 | -2 | 0)$; Deckkante $E(1 | 3 | 2)$, $F(1 | -1 | 2)$. Zeige: $ABC$ ist gleichschenklig. Begründe: $EF$ ist parallel zur Grundfläche. \hfill (Abitur 2023 GK)
\teil $D(8 | 0 | 2)$, $E(0 | 6 | 2)$, $F(0 | 0 | 2)$ und $M(4 | 3 | 2)$. Begründe: $D$, $E$ und $F$ liegen auf einem Kreis mit dem Mittelpunkt $M$. \hfill (Abitur 2024 GK)
\teil In der $x_1x_2$-Ebene liegen $A(0 | 0)$, $B(6 | 0)$ und $C(3 | 9)$ auf einem Kreis. Koordinaten seines Mittelpunkts $M$? \hfill (Abitur 2019 GK) \\ M( \leerfeld | \leerfeld )
\teil $A(0 | 6 | 1)$, $B(0 | -6 | 1)$ und $C(0 | 0 | p)$ mit $p > 1$. Der Umfang des Dreiecks $ABC$ ist $32$. $p$? \hfill (Abitur 2024 LK) \\ p = \leerfeld
\teil Dreieck $OPQ$ mit dem Ursprung $O$, $P(4 | 2 | 1)$ und $Q(1 | 5 | 3)$. Gib einen Punkt $T \ne P$ an, sodass das Dreieck $OTQ$ denselben Flächeninhalt und denselben Umfang hat. \hfill (Abitur 2023 LK) \\ T( \leerfeld | \leerfeld | \leerfeld )
\teil Dreieck $A(2 | 1 | 2)$, $B(4 | 5 | 4)$, $C(5 | 3 | 1)$ mit $|CA| = |CB|$. Gesucht ist ein Punkt $D \ne C$ mit $|DA| = |DB|$ und $|DA| = |CA|$. Koordinaten von $D$? \hfill (Abitur 2023 GK) \\ D( \leerfeld | \leerfeld | \leerfeld )
\end{teile}
\end{aufgabe}

%% TODO Titel: Ich-kann-Satz fehlt in der Bank (Platzhalter: Kettenname) – Nr. 16
%% TODO Anweisung: ein Satz über der ersten Teilaufgabe fehlt in der Bank – Nr. 16
\begin{aufgabe}{Dreiecke nachweisen – Fehler finden, Begründen}
\begin{teile}
\teil Ich finde den Fehler: Zu zeigen ist, dass $A(2 | 3 | 1)$, $B(6 | 3 | 1)$, $C(4 | 5 | 2)$ ein gleichschenkliges Dreieck mit der Basis $AB$ bilden. Tom: „$|AB| = 4$, $|AC| = 3$ – nicht gleichschenklig.“
\teil Begründe: Für „gleichschenklig“ genügt es, zwei Seitenlängen zu vergleichen. Welche?
\end{teile}
\end{aufgabe}
